\documentclass[a4paper, 12pt]{article}
\usepackage{kotex}

\usepackage[a4paper, margin=2.5cm]{geometry}
\usepackage{tocloft}
% Preamble에 추가
\usepackage{microtype}
\usepackage{setspace}
\usepackage{tikz} 
\usepackage{graphicx}
\usepackage{float} % [H] 옵션을 사용하기 위해 반드시 추가해야 합니다.
\usepackage{booktabs}   % 고품질 표 작성을 위한 패키지
\usepackage{tabularx}   % 표 너비 자동 조절


\begin{document}

% 1. 타이틀 전용 페이지 시작 %%%%%%%%%%%%%%%%%%%%%%%%%%%%%%%%%%%%
\begin{titlepage}
    \centering
    \vspace*{3cm}
    {\large \textbf{청덕고등학교 3학년 과학과제연구}} \par
    \vspace*{2.5cm}
    
    {\Huge \textbf{멀티 프로세서와 단일 프로세서의} \par}
    {\Huge \textbf{작업 효율성 비교분석} \par}
    \vspace{10.0cm}
    {\Large \textbf{TES}\par}
    \vspace{2cm}
    {\large 청덕고등학교 \par}
    
    \vfill % 나머지 공간을 채움
    
    {\large 작성시작일: 2026-3-14 \par}
    \vspace{2cm}
\end{titlepage}
%%%%%%%%%%%%%%%%%%%%%%%%%%%%%%%%%%%%%%%%%%%%%%%%%%%%%%%%%
\newpage
% 1. 기본 제목을 비워버립니다. (정렬 문제 해결)
\renewcommand{\contentsname}{}

% 2. 수동으로 가운데 정렬된 제목을 만듭니다.
\begin{center}
    \Large\bfseries 차\kern 30pt 례
\end{center}
\tableofcontents

\newpage % 차례가 끝나고 다음 페이지에서 본문을 시작하게 함
%%%%%%%%%%%%%%%%%%%%%%%%%%%%%%%%%%%%%%%%%%%%%%%%%%%%%%%%%
\section{연구의 요약}
\subsection{연구동기, 목적}
\begin{spacing}{1.5}
\fontdimen2\font=0.7em
컴퓨터에 사용되는 일부 프로그램등은 컴퓨터의 CPU에 큰 부담을 주기도 하는데 이때 시스템 모니터링 기능을 통해 CPU의 사용량을 확인하면 CPU의 코어 일부에만 사용량이 집중되어 있기도 한다. \par
이러한 경험을 통해 CPU의 병렬화, 부하 분산에대한 관심을 가지게 되었으며 이에 '작업의 분산의 효율성'에 초점을 맞추어 작업을 분산하는것이 그렇지 않은것보다 정확히 어떤 면에서 장점을 가지는지에 대해 연구를 계획하였다.
\end{spacing}
%%%%%%%%%%%%%%%%%%%%%%%%%%%%%%%%%%%%%%%%%%%%%%%%%%%%%%%%%%%%%%%%%%%%%%%%%%%%%%%%%%%%%%%%%%
%\newpage
\section{연구의 실행}
\subsection{이론적 배경 및 선행연구}
\begin{spacing}{1.5}
\fontdimen2\font=0.7em
컴퓨터 프로세서(Processor)란 컴퓨터 시스템의 핵심 부품으로, 외부로부터 입력받은 데이터를 프로그램의 명령에 따라 처리, 연산하고 그 결과를 출력하는 중앙 처리 장치(CPU)를 의미한다. 멀티 프로세서는 하나의 컴퓨터 시스템 내에 2개 이상의 중앙처리장치(CPU)를 장착하여, 메모리와 입출력 장치를 공유하며 병렬 처리를 수행하는 시스템을 말한다. 

멀티프로세서는 크게 대칭형 멀티프로세서읜 SMP와 비대칭형 멀티프로세서인 AMP 분류되며 SMP의 경우 모든 프로세서가 운영체제나 작업 관리 면에서 동일한 권한을 가지며 이는 어떤 프로세서든 시스템 내의 모든 메모리와 I/O에 접근할 수 있어 부하 분산에 유리함을 뜻한다. AMP의 경우는 주(Master) 프로세서가 시스템을 제어하고 스케줄링을 담당하며, 부(Slave) 프로세서는 주 프로세서의 명령에 따라 특정 작업만 수행한다.
\end{spacing}

\subsection{연구주제의 선정}
\begin{spacing}{1.5}
\fontdimen2\font=0.7em
본 연구에서는 컴퓨터 시스템의 핵심인 프로세서의 병렬화가 작업 효율성에 미치는 영향을 실험적으로 분석하기 위해 '멀티 프로세서와 단일 프로세서의 작업 효율성 비교분석'을 주제로 선정하였다. 이는 이론적인 개념을 넘어 실제 하드웨어를 구축하여 직접 측정함으로써, 병렬 처리의 장점을 체감할 수 있는 기회를 제공한다. 특히, Arduino 플랫폼을 활용한 간단한 멀티 프로세서 시스템을 통해 교육적 가치가 높으며, 현대 컴퓨터 아키텍처의 기초를 이해하는 데 도움이 된다.
\end{spacing}

\subsection{연구 방법 및 과정}
\begin{spacing}{1.5}
\fontdimen2\font=0.7em
본 연구는 ATmega 328P 마이크로 컨트롤러와 Arduino Nano 보드를 기반으로 한 멀티 프로세서 시스템을 구축하여 수행되었다. 먼저, 두 개의 ATmega328P 마이크로컨트롤러를 프로세서로 사용하여 공유 메모리(RAM)를 통해 통신하는 시스템을 설계하였다. 하드웨어로는 74HC595 시프트 레지스터, 74HC245 버퍼, 62256 RAM 등을 사용하여 주소 및 데이터 버스를 구성하였다.

전체적인 연구 과정은 다음과 같다:
\begin{enumerate}
    \item {하드웨어 설계 및 구축:} 각 하드웨어들의 데이터 시트를 통해 핀 연결 목록을 PinManual.md에 명시하여 회로를 구성하고, 검증 및 디버깅을 통해 시스템을 테스트하고 핀 연결 목록을 수정하였다.
    \item {소프트웨어 개발:} SMU(System Memory Unit) 프로그램을 통해 RAM에 프로그램을 업로드하고, 프로세서들이 병렬로 실행할 수 있도록 하였다. 따라서 프로세서에서는 프로세서가 스스로 RAM에 접근하여 데이터를 읽어올수 있고 이를 해석하여 동작하는 코드를 업로드 하였다. Python 기반의 사용자 인터페이스(mainMPE.py)를 개발하여 편리한 프로그래밍을 지원하였다.
    \item {실험 설계:} 동일한 계산 작업(예: 반복적인 산술 연산)을 단일 프로세서와 멀티 프로세서 환경에서 수행하도록 프로그램을 작성하였다. 일정시간 동안 처리한 명령어의 개수를 비교를 통해 효율성을 비교하였다.
    \item {데이터 수집 및 분석:} 다른종류의 코드를 실행시켜보며 전체적이고 평균적인 환경에서 병렬화의 이점을 분석하였다.
\end{enumerate}
\end{spacing}
%%%%%%%%%%%%%%%%%%%%%%%%%%%%%%%%%%%%%%%%%%%%%%%%%%%%%%%%%%%%%%%%%%%%%%%%%%%%%%%%%%%%%%%%%
\subsubsection{시스템 동작 알고리즘 설계}
% 대문자 H를 사용하여 위치를 강제 고정합니다.
\begin{figure}[H] 
    \centering
    \includegraphics[width=1\linewidth]{corearch.png}
    \caption{Main Architecture}
    \label{fig:placeholder}
\end{figure}
\begin{spacing}{1.5}
\fontdimen2\font=0.7em
부하의 분산이라는 개념에 중점을 두어 설계한 본 연구의 듀얼 프로세서의 기본 구조는 Figure 1과 같다.

시스템의 데이터 버스의 크기는 8비트 이며 주소 버스의 크기는 RAM의 용량에 따라 변한다. 본 연구에서는 256Kb(32KB) 크기의 62256 SRAM을 이용하였으며 주소버스의 크기가 14비트이기 때문에 데이터버스 8핀 + 주소버스 14핀 + 하드웨어 제어에 필요한 나머지 핀까지 해서 프로세서에는 총 22개이상의 입출력이 가능한 핀이 있어야 한다. 

%%%%%%%%%%%%%%%%%%%%%%%%%%%%%%%%%%%%%%%%%%%%%%%%%%%%%%%%%%%%%%%%%
\begin{figure}
    \centering
    \includegraphics[width=1\linewidth]{atmega328pdip.png}
    \caption{ATmega328P DIP 패키지 핀 맵(ATMEGA328P : 8-bit Microcontroller with 4/8/16/32K Bytes In-System Programmable Flash
ATMEL Corporation)}
    \label{fig:placeholder}
\end{figure}
%%%%%%%%%%%%%%%%%%%%%%%%%%%%%%%%%%%%%%%%%%%%%%%%%%%%%%%%%%%%%%%
그러나 본 연구에서의 프로세서는 ATmega328P MCU이며 RESET 핀과 XTAL1,2 핀을 제외하면 PD0\textasciitilde7, PB0\textasciitilde5, PC0\textasciitilde5 까지 총 20개로 RAM과의 주소선을 전부 연결하기에는 핀의 개수가 부족하다.

이를 극복하기 위해 \textbf{74HC595}를 이용하였다. 
74HC595는 8비트 시프트 레지스터로, 시리얼 데이터를 병렬로 출력하는 IC이다. 이 칩은 SER(시리얼 입력), SCK(시프트 클럭), RCK(래치 클럭), OE(출력 인에이블) 핀을 통해 제어된다. 

동작 방식은 다음과 같다: 시프트 클럭(SCK) 펄스마다 SER 핀의 데이터를 내부 레지스터로 시프트시키고, 래치 클럭(RCK) 펄스가 주어지면 내부 레지스터의 데이터를 출력 핀(QA\textasciitilde QH)으로 전송한다. OE 핀이 로우일 때 출력이 활성화된다.


본 시스템에서 프로세서는 각각 한 개의 74HC595를 연결하여 주소 버스의 상위 비트(A8\textasciitilde A13)를 제어하며 SER 핀을 통해 데이터를 순차적으로 전송한다.
A0\textasciitilde A7는 ATmega328P의 핀들이 직접 담당하며, 
이와같은 설계의 의도는 후술할 RAM의 뱅크 스위칭와 페이징 방식 설계에 의한 것이다.
이를 통해 ATmega328P의 제한된 핀 수로도 14비트 주소 버스를 완전히 제어할 수 있다.
\end{spacing}
%%%%%%%%%%%%%%%%%%%%%%%%%%%%%%%%%%%%%%%%%%%%%%%%%%%%%%%%%%%%%%%%%%%%%%%%%%%%%%%%%%%%%%%%%%%%%%%%%
\begin{spacing}{1.5}
\begin{enumerate}
\fontdimen2\font=0.7em
\item \textbf{뱅크 스위칭과 페이징 시스템} 

두개의 프로세서(MCU)가 공유 메모리에 접근해서 RAM내에서 충돌없이 프로세서가 각자의 작업을 수행하기 위해서는 RAM 내부의 영역을 프로세서1과 프로세서2의 영역 두개로 나누어야 한다고 판단하였다.

이를 이용해 RAM의 A14핀을 74HC595를 사용해 절약한 프로세서1의 남는 핀()에 연결하고 프로세서1이 RAM으로부터 데이터를 읽고난 후 해당 핀을 활성화하여 프로세서2가 사용할 RAM의 영역 다른 절반 부분을 활성화 시키도록 뱅크 스위칭 방식을 설계하였다. 프로세서2는 RAM으로부터 데이터를 읽고난 후 \textbf{작업완료 알림용  신호 핀}을 통해 프로세서1에 알리고 프로세서1은 다시 뱅크를 자신의 영역으로 활성화 한다. \par

페이징은 메모리를 고정된 크기의 블록으로 나누어 관리하는 메모리 관리 기법이며 프로세서가 접근할수 있는 메모리(RAM)의 용량보다 실제 장착된 메모리의 크기가 더 크고, 그럼에도 메모리 전체를 활용해야 하기에 이를 도입하였다.

프로세서의 상위 7비트(A0\textasciitilde A6 고정된 128바이트 규격)는 프로세서의 핀에 직접 연결하고, 나머지 하위 비트(A7\textasciitilde A13: 74HC595를 이용해 제어하는 영역)은 페이징 시스템에 이용된다.

이와 같다면 프로세서가 모든 RAM의 주소에 페이징 없이 접근할수 있는것처럼 확인되지만
74HC595는 직렬을 병렬로 전송해주는 로직이며 병렬로 입력을 받은것을 다시 직렬로 변환하는 기능은 없기 때문에 74HC595로 RAM의 주소선을 건드려 상위 7비트로만 데이터를 받아들이는 것이다.

프로세서가 RAM의 500번지 \((500 = 128 \times 3 + 116)\)에 있는 데이터를 가져오려 할 때, 
본 연구의 페이징 시스템은 다음과 같이 작동된다.

\begin{itemize}
    \item \textbf{1. 페이지 번호 계산 :}
    프로세서는 접근하려는 주소를 128로 나눕니다. 몫(3)은 페이지 번호가 되고, 나머지(116)는 페이지 내 오프셋이 된다.
    
    \item \textbf{2. 페이지 스위칭 (페이지 선택 단계) :}
    프로세서는 74HC595에 3번 페이지를 로드한다. PD0(SER), PD1(SCK), PC3(RCK)를 사용하여 74HC595에 '3'이라는 값을 넣는다. 
    이제 메모리의 $A_7 \sim A_{13}$ 핀이 '3'으로 고정된다. 즉, "3번 페이지 구역"이 활성화되었다는 의미이다.

    \item \textbf{3. 오프셋 주소 적용 (데이터 접근 단계) :}
    프로세서는 주소 버퍼(74HC245)를 통해 하위 7비트(116)를 메모리의 $A_0 \sim A_6$ 핀으로 보낸다.
\end{itemize}

\item  \textbf{데이터 버스와 주소버스의 작동구조}

Figure 1에서 알수 있듯 두 프로세서는 같은 주소버스와 데이터 버스를 공유한다. 때문에 두개의 프로세서는 버스에서의 충돌로 인해 동시에 RAM에 접근할수 없으며 두 프로세서가 번갈아가며 데이터, 주소 버스를 점유하는 방식으로 동작하도록 설계하였다. 

각 프로세서는 \textbf{74HC245}라는 버스 트랜시버 칩을 가지게 되는데 이것이 버퍼의 역할을 하여 각 프로세서의 내부 데이터/주소 버스를 시스템 공통 버스로부터 전기적으로 격리(Isolation)하는 게이트웨이가 된다.

평상시 프로세서1이 버스를 점유하고 있을 때는 프로세서1의 74HC245 버퍼가 '활성(Enable)' 상태를 유지하여 메모리와 신호를 주고받지만, 이때 프로세서2의 버퍼는 '하이 임피던스(High-Impedance, Hi-Z)' 상태가 되어 시스템 버스로부터 완전히 차단된다. 이러한 구조는 버스 상의 전기적 충돌(Short-circuit)을 원천적으로 차단하며, 특정 프로세서가 버퍼를 비활성화하는 즉시 해당 프로세서는 메모리 버스에서 논리적으로 제거되는 효과를 만든다.

특히, 두 프로세서의 버퍼 제어 신호는 \textbf{74HC04(Hex Inverter)}를 통해 상호 배타적(Mutual Exclusive)으로 동작하도록 설계되었다. 74HC04는 입력된 논리 레벨을 반전시켜 출력하는 소자로, 본 시스템에서는 하나의 제어 신호를 입력받아 프로세서1의 버퍼가 활성화될 때 프로세서2의 버퍼는 반드시 비활성화되도록 물리적인 논리 경로를 강제한다. 이는 소프트웨어적 오류로 인해 두 프로세서가 동시에 버퍼를 열어버리는 '버스 충돌'을 하드웨어 레벨에서 원천적으로 차단하기 위한 필수적인 안전장치이다.

\item  \textbf{버스 소유권 전환 방법}

\textbf{1. 뱅크 스위칭과 페이징 시스템}에서 사용한 것과 같은 \textbf{작업완료 알림용 신호 핀}을 적극 사용하여, 프로세서1이 프로세서2에 버스 접근 권한을 양도하는 방식으로 동작한다. 74HC04를 통한 인버팅 로직을 거치기 때문에, 한 프로세서의 버퍼가 'Low' 신호를 받아 열리면(OE 활성화), 반대편 프로세서의 버퍼는 인버터를 거친 'High' 신호를 받아 즉시 닫히게(OE 비활성화) 된다. 구체적인 전환 절차는 다음과 같다.
\begin{itemize}
\item \textbf{1단계 (점유 해제):} 프로세서1은 현재 수행 중인 메모리 작업을 완료한 후, 74HC04에 입력되는 제어 신호를 토글하여 자신의 74HC245 버퍼를 하이 임피던스(Hi-Z) 상태로 전환한다.
\item \textbf{2단계 (권한 이양):} 프로세서1은 \textbf{작업완료 알림용 신호 핀}을 통해 프로세서2에 버스 사용이 가능함을 알린다. 이 신호를 받은 프로세서2는 자신의 제어 신호를 활성화하여 74HC245 버퍼를 엶으로써 시스템 공통 버스의 제어권을 획득한다.
\item \textbf{3단계 (데이터 접근):} 프로세서2는 자신의 페이지를 활성화하고, 페이징 시스템을 통해 메모리의 특정 주소에 접근한다. 이때 프로세서1의 모든 출력 핀은 74HC04에 의해 전기적으로 차단된 상태이므로 데이터 왜곡이 발생하지 않는다.
\item \textbf{4단계 (복귀):} 프로세서2는 작업이 완료되면 다시 알림 신호를 프로세서1로 보내고, 버퍼를 비활성화하여 버스 점유권을 반환한다.
\end{itemize}

\end{enumerate}
\end{spacing}
%%%%%%%%%%%%%%%%%%%%%%%%%%%%%%%%%%%%%%%%%%%%%%%%%%%%%%%%%%%%%%%%%%%%%%%%%%%%%%%%
\subsubsection{프로세서 코드}
\begin{spacing}{1.5}
\fontdimen2\font=0.7em    
프로세서는 8비트의 데이터 크기중 상위 4비트를 동작할 명령어 + 하위 4비트를 즉시값으로 취급하며 수행할수 있는 명령어의 종류는 트랜지스터를 이용한 게이트 로직, 간단한 어셈블리어 명령어들로 구성하였다. \par
\vspace{1cm}
\begin{table}[H]
    \centering
    \caption{Instruction Set  Definitions}
    \vspace{0.8em}
    
    \begin{tabularx}{\textwidth}{l l l X}
        \toprule
        \textbf{Mnemonic} & \textbf{Opcode} & \textbf{Operand} & \textbf{Description} \\
        \midrule
        \texttt{NOP}      & \texttt{0x00} & - & 아무 동작도 수행하지 않음 (No Operation). \\
        \texttt{LOAD}     & \texttt{0x10} & Imm & \texttt{레지스터 A}에 상숫값(4-bit)을 직접 로드함. \\
        \texttt{ADD}      & \texttt{0x20} & Imm & \texttt{regA}에 상숫값을 더함. \\
        \texttt{SUB}      & \texttt{0x30} & Imm & \texttt{regA}에서 상숫값을 뺌. \\
        \texttt{MUL}      & \texttt{0x40} & Imm & \texttt{regA}에 상숫값을 곱함. \\
        \texttt{AND}      & \texttt{0x50} & Imm & \texttt{regA}와 상숫값 간 비트 AND 연산 수행. \\
        \texttt{OR}       & \texttt{0x60} & Imm & \texttt{regA}와 상숫값 간 비트 OR 연산 수행. \\
        \addlinespace
        \texttt{OUT}      & \texttt{0x70} & - & \texttt{regA}의 값을 출력 레지스터로 전송. \\
        \texttt{FETCH}    & \texttt{0x80} & Imm & 지정된 슬롯에서 값을 읽어 \texttt{regA}에 로드. \\
        \texttt{SLOT}     & \texttt{0x90} & Imm & \texttt{regA}의 값을 지정된 메모리 슬롯에 저장. \\
        \addlinespace
        \texttt{PUSH}     & \texttt{0xA0} & - & \texttt{regA}의 값을 스택에 저장 (Max: 8-depth). \\
        \texttt{POP}      & \texttt{0xB0} & - & 스택 상단의 값을 꺼내어 \texttt{regA}에 로드. \\
        \texttt{SETPAGE}  & \texttt{0xC0} & - & \texttt{74HC595}를 통해 페이지 전환 및 \texttt{PC} 초기화. \\
        \texttt{HALT}     & \texttt{0xD0} & - & 프로세서 동작 정지. \\
        \bottomrule
    \end{tabularx}
\end{table}

4비트 명령어와 4비트 즉시값(Immediate)으로 구성된 명령어 구조로 인해, 
\texttt{LOAD} 명령어를 통해 레지스터에 즉시 로드할 수 있는 최댓값은 15(0~15)로 제한되며, 
이는 \texttt{ADD}와 같은 연산 명령어에서도 동일한 제약으로 작용한다. 
이러한 하드웨어적 한계를 극복하기 위해 본 시스템은 \textbf{스택(Stack)}과 \textbf{변수(Slot)} 체계를 도입하였다. 
기존 명령어들의 연산 결괏값을 \textbf{PUSH}하면 해당 값은 스택에 저장되어 추후 \textbf{POP}을 통해 다시 참조할 수 있다.

변수 체계는 \textbf{SLOT}으로 명명하였으며, 0번부터 15번까지 총 16개의 슬롯을 제공한다. 
명령어는 4비트 명령어와 4비트 주소값(Slot Index)으로 구성되지만, 
실제 슬롯 내부에 저장되는 데이터는 8비트 형태이므로, 
프로세서는 슬롯을 활용하여 4비트 즉시값 제약을 넘어서는 8비트 크기의 데이터를 자유롭게 연산할 수 있다.

\end{spacing}

\section{연구 결과 및 시사점}
\begin{spacing}{1.5}
\fontdimen2\font=0.7em
시사점으로는, 멀티 프로세서의 도입이 시스템의 성능을 극대화할 수 있음을 확인하였다. 그러나, 병렬화에는 추가적인 하드웨어 비용과 복잡성이 따르므로, 작업의 성격에 따라 적절한 선택이 필요하다. 본 연구를 통해 컴퓨터 아키텍처의 기초를 이해하고, 미래의 고성능 컴퓨팅에 대한 관심을 높일 수 있었다.

그 과정에서 배운 것들: 기술적/학술적 교훈
AVR GPIO의 창의적 활용: AVR의 디지털 I/O를 데이터/주소 버스로 확장해 외부 메모리를 액세스하는 법을 배웠어요. PD2PD7을 데이터 버스로, PB2PC2를 주소 버스로 사용하며, 타이밍 동기화(SYNC(언더바)DELAY 매크로)를 통해 신호 지터를 최소화. 이는 "하드웨어 제약을 소프트웨어로 해결"하는 방법을 가르쳤습니다.
멀티프로세서 동기화의 복잡성: PC4 핀으로 버스 마스터리를 제어하고, PC5로 인터럽트 트리거를 전송하며 듀얼코어를 조율하는 경험. 경쟁 조건(Race Condition)을 방지하기 위해 핸드셰이크 프로토콜을 구현한 점에서, 분산 시스템의 타이밍과 신호 무결성을 배웠어요.

VMF 설계와 최적화: 명령어 세트(NOP, LOAD, ADD 등)를 정의하고, 버스트 모드(4개 명령어 연속 실행)로 오버헤드를 75퍼센트로 줄인 것. AVR의 느린 GPIO를 극복하기 위해 delay 제거와 캐싱(페이지 캐시)을 적용 – 이는 컴퓨터 아키텍처의 ISA(Instruction Set Architecture)와 메모리 계층화를 실전에서 체득한 교훈입니다.
디버깅과 검증의 중요성: 로직 애널라이저로 엣지 타이밍을 측정하고, 멀티미터로 전압 강하를 확인하며 시스템을 튜닝. STAGE별 프로토타이핑으로 "작은 실패를 쌓아 큰 성공"을 이루는 법을 배웠어요. 이는 전자공학에서 "증분 설계"의 실천적 가르침입니다.

3. 이 연구의 의의: 학술적/기술적 기여
창의적 문제 해결의 모델: AVR의 제약을 인정하면서도, 멀티프로세서 시스템을 구현한 것은 "제약이 혁신의 촉매"라는 교훈을 줍니다. 이는 전자공학 연구에서 "비표준 플랫폼으로 표준 문제를 해결"하는 사례로, FPGA나 ASIC 없이 MCU로 복잡한 시스템을 설계하는 방법을 제시합니다. proc2VMF.ino의 최적화 효과(듀얼코어로 8배 성능 향상)는 AVR의 잠재력을 극대화한 증거입니다.

교육적 가치의 확대: 이 프로젝트는 AVR의 한계를 극복한 실험으로, 학생/연구자에게 "왜 AVR이 멀티프로세서용이 아닌지"를 가르치면서도, 그것을 가능하게 하는 기술(인터페이스 설계, 타이밍 최적화)을 전수합니다. 로직 애널라이저 기반 디버깅을 강조하여, 실전 전자공학 스킬을 키웁니다.(?)

기술적 혁신의 기여: VMF를 통해 가상 머신을 구현한 점은 임베디드 시스템에서 "소프트웨어 정의 하드웨어"의 가능성을 입증. 멀티코어 시뮬레이션으로, 저비용 플랫폼에서 고급 컴퓨팅 개념(멀티태스킹, 메모리 관리)을 탐구하게 합니다. 이는 IoT나 임베디드 연구에서 AVR을 "실험 도구"로 재정의하는 의의를 가집니다.

미래 연구의 기반: 이 AVR 기반 멀티프로세서는 Z80 교체의 프로토타입으로, "MCU 한계를 넘어선 시스템 설계"의 연구 방향을 제시합니다. 보람은 단순한 성공이 아니라, "배운 것을 적용해 더 큰 문제를 해결"하는 동기부여입니다.

\end{spacing}









\section{참고 문헌}
\begin{spacing}{1.5}
\fontdimen2\font=0.7em
\begin{enumerate}
    \item Arduino 공식 문서
    \item ATmega328P 데이터시트: Microchip Technology Inc.
\end{enumerate}
\end{spacing}

\newpage{table}

\begin{table}[H] % [H]는 '정확히 여기(Here)'에 고정하라는 강력한 명령입니다.
    \centering
    \caption{Hardware Components List}
    \vspace{0.8em}
    
    % tabularx를 사용하여 표 전체 너비를 문서 너비(\textwidth)에 맞춥니다.
    % L, C, R은 각각 왼쪽, 중앙, 오른쪽 정렬 기능을 가진 가변 너비 열입니다.
    \begin{tabularx}{\textwidth}{l X c}
        \toprule
        \textbf{parts type} & \textbf{name} & \textbf{number} \\
        \midrule
        MCU & \texttt{Arduino Nano} & 2 \\
        MCU & \texttt{Atmega328P (DIP)} & 2 \\
        Memory & \texttt{62256 SRAM} & 1 \\
        \addlinespace
        IC logic & \texttt{74HC245} & 4 \\
        IC logic & \texttt{74HC595} & 4 \\
        IC logic & \texttt{74HC04} & 1 \\
        IC logic & \texttt{74HC125} & 1 \\
        IC logic & \texttt{74HC165} & 2 \\
        \addlinespace
        Oscillator & \texttt{Active Crystal} & 1 \\
        \bottomrule
    \end{tabularx}
\end{table}




\end{document}